\section*{第 \ref{ch:b2:mammal-reproduction} 章——哺乳动物的有性生殖}

\begin{solution}{exo:b2:mammal-reproduction:1}
精原细胞自青春期起终生分裂；卵原细胞
仅在胎儿期分裂。一次\omterm{def:b2:meiosis-heredity:meiosis}{减数分裂}产生四个精子，但只有一个卵子
（另有三个极体）。精子形成约需 64 天，外加两周的成熟期；
卵子自出生前即停滞于前期 I，直至其\omterm{def:b2:mammal-reproduction:ovary}{排卵}之月——
历时数十年——之后又停滞于中期 II，直至受精。
男性每天产生 $10^{8}$ 个精子，一生约 $10^{12}$ 个；
女性一生约\omterm{def:b2:mammal-reproduction:ovary}{排卵} 400 个。
\end{solution}

\begin{solution}{exo:b2:mammal-reproduction:2}
与\omterm{prop:b2:mammal-reproduction:fertilisation}{透明带}结合；\omterm{prop:b2:mammal-reproduction:fertilisation}{顶体反应}，消化出一条
通道；细胞膜融合，精子细胞核进入；
钙离子波与皮层反应（\omterm{prop:b2:mammal-reproduction:fertilisation}{透明带}变硬，失去
其受体，阻止多精入卵）；完成\omterm{def:b2:meiosis-heredity:meiosis}{减数分裂} II，
排出第二极体；形成两个原核，
DNA 复制，进行第一次有丝分裂。
\end{solution}

\begin{solution}{exo:b2:mammal-reproduction:3}
\omterm{def:b2:mammal-reproduction:implantation}{滋养层}（外层）形成\omterm{def:b2:mammal-reproduction:placenta}{胎盘}与胎膜；\omterm{def:b2:mammal-reproduction:implantation}{内细胞团}
发育为胚胎。\omterm{def:b2:mammal-reproduction:implantation}{着床}始于受精后第 6 至 7 天，
在由孕酮准备好的子宫内膜中进行，
孕酮来自\omterm{def:b2:mammal-reproduction:ovary}{黄体}。
\end{solution}

\begin{solution}{exo:b2:mammal-reproduction:4}
氧气：扩散。葡萄糖：易化运输。母体红细胞
不能通过（母婴血液从不混合）。IgG：受体介导的
转运。尿素：由胎儿向母体扩散。酒精：
可自由扩散。细菌：多数不能通过（少数如
梅毒和李斯特菌等可以）。
\end{solution}

\begin{solution}{exo:b2:mammal-reproduction:5}
$m = 3, 1, 0.4, 0.1$：$P = 0.95, 0.63, 0.33, 0.095$。
概率低于二分之一的条件是 $m < \ln 2 = 0.69$，即低于 $7\times 10^{7}$ 个精子。
曲线趋于饱和：超过约 $10^{8}$ 个额外精子后几乎
无作用；低于此值时，生育力几乎与精子数量成正比地下降。
\end{solution}

\begin{solution}{exo:b2:mammal-reproduction:6}
青春期前，13 年（\num{4750} 天）中损失 $7\times 10^{5}$ 个，
约每天 150 个。青春期后，37 年（\num{13500}
天）中损失 \num{299000} 个，即每天 22 个；\omterm{def:b2:mammal-reproduction:ovary}{排卵}为 $400/\num{13500} = 0.03$ 个/天。闭锁
\omterm{def:b2:mammal-reproduction:ovary}{卵泡}数为\omterm{def:b2:mammal-reproduction:ovary}{排卵}数的七百倍；卵子储备的消耗
并非源于\omterm{def:b2:mammal-reproduction:ovary}{排卵}，而是源于死亡。
\end{solution}

\begin{solution}{exo:b2:mammal-reproduction:7}
成体：$(30/27)^{2.7} = 1.33$，$S = 1.33/2.33 = 0.57$。胎儿：
$(30/19)^{2.7} = 3.4$，$S = 3.4/4.4 = 0.77$。在\omterm{def:b2:mammal-reproduction:placenta}{胎盘}的低氧分压下，
胎儿血液比成体血液多载约五分之一的氧，
并在组织中卸氧，此时成体血红蛋白
已有一半空载。
\end{solution}

\begin{solution}{exo:b2:mammal-reproduction:8}
$2^{k} = 25$：$k = 4.6$ 次倍增，\qty{9.3}{d}：受精后第 16 至 17 天，
即末次月经后第 30 至 31 天——约为
漏经当日。
\end{solution}

\begin{solution}{exo:b2:mammal-reproduction:9}
$3\times 10^{-8}\times 6\times 10^{4}\times 20/10^{-3} =
\qty{36}{mL/min}$，而需求为 \qty{25}{mL/min}：余量
由 8 倍降至 1.5 倍；胎儿在任何额外负荷下
都会缺氧，生长迟缓，可能需提前分娩。
\end{solution}

\begin{solution}{exo:b2:mammal-reproduction:10}
(a) \omterm{def:b2:mammal-reproduction:testis}{睾丸}形成并分泌两种激素；AMH 使苗勒管
退化，但睾酮无法发挥作用：沃尔夫管
退化，外生殖器为女性型——即有\omterm{def:b2:mammal-reproduction:testis}{睾丸}而无
子宫的女性。(b) 睾酮发挥作用，AMH 不起作用：
男性，有附睾、输精管和男性外生殖器，另有子宫和输卵管。(c)
卵巢及苗勒管衍生结构（无\omterm{def:b2:mammal-reproduction:testis}{睾丸}、无 AMH）；肾上腺
雄激素使外生殖器不同程度男性化。(d)
\emph{\omterm{prop:b2:mammal-reproduction:sex}{SRY}} 促使形成\omterm{def:b2:mammal-reproduction:testis}{睾丸}，\omterm{def:b2:mammal-reproduction:testis}{睾丸}引导男性发育途径：一名具有两条
X 染色体的男性，因\omterm{def:b2:mammal-reproduction:testis}{精子发生}所需的 Y 基因缺失
而不育。
\end{solution}

\begin{solution}{exo:b2:mammal-reproduction:11}
最常见的（\qty{68}{\%}）为第 4 至第 8 天发生的分裂：在
\omterm{def:b2:mammal-reproduction:implantation}{滋养层}已形成（因此只有一个\omterm{def:b2:mammal-reproduction:placenta}{胎盘}）之后、羊膜
形成之前（因此有两个囊）。第 3 天之前，每一半各自形成\omterm{def:b2:mammal-reproduction:implantation}{滋养层}，
因而形成两个\omterm{def:b2:mammal-reproduction:placenta}{胎盘}（\qty{30}{\%}）；第 8 天之后，
羊膜已形成并被共享（\qty{2}{\%}）。\omterm{def:b2:mammal-reproduction:implantation}{内细胞团}
在\omterm{def:b2:mammal-reproduction:implantation}{滋养层}定型之后才分裂，显然是
最容易发生的事件。
\end{solution}

\begin{solution}{exo:b2:mammal-reproduction:12}
有袋类在出生前投入很少：妊娠期短，幼体极小，
哺乳期可以缩短——若遇干旱，母体会丢弃育儿袋中的幼体，
因此每次繁殖的风险较低，母体可立即再次尝试。有\omterm{def:b2:mammal-reproduction:placenta}{胎盘}类
则大量投入于较长的妊娠期，产下大而发育良好、出生时
存活率更高的幼体，代价是母体数月受制，
期间较为脆弱。在产出丰富且可预测的环境中，
\omterm{def:b2:mammal-reproduction:placenta}{胎盘}类幼体较高的存活率占优势；澳大利亚气候多变，
有利于有袋类放弃幼体的选择；其隔离又
使\omterm{def:b2:mammal-reproduction:placenta}{胎盘}类长期未能进入，
直至近代。
\end{solution}

\begin{solution}{pb:b2:mammal-reproduction:1}
\textbf{1.} $m = 3\times 10^{8}\times 10^{-8} = 3$; $P = 1 - e^{-3} =
0.95$.
\textbf{2.} $1 - e^{-3}(1 + 3) = 0.80$：若无阻断，
则五次受精中有四次为\omterm{thm:b2:mammal-reproduction:poisson}{多精受精}，胚胎为三倍体。
\textbf{3.} $0.15/5\times 10^{-5} = \qty{3000}{s}$，约 50 分钟；
由子宫与输卵管的收缩将其输送。
\textbf{4.} $300/3\times 10^{8} = 10^{-6}$：在酸性阴道中丢失，
在宫颈黏液中受阻，被子宫内的吞噬细胞清除，以及
进入错误的输卵管。
\textbf{5.} \omterm{def:b2:mammal-reproduction:ovary}{排卵}前三天至\omterm{def:b2:mammal-reproduction:ovary}{排卵}后一天：约
4 天。
\textbf{6.} $m = 0.5$，$P = 0.39$；期望周期数 $1/P = 2.5$。
\textbf{7.} 精子中段的少量线粒体被标记，
并在入卵后被清除；卵子则携带十万个。
\textbf{8.} $120/16 = 7.5$：7 次分裂，$2^{7} = 128$ 个细胞——
约一百个。
\textbf{9.} $\log_2 25 = 4.6$ 次倍增，\qty{9.3}{d}：第 16 至 17 天；
即末次月经后第 30 至 31 天。
\textbf{10.} \omterm{def:b2:mammal-reproduction:ovary}{黄体}在\omterm{def:b2:mammal-reproduction:ovary}{排卵}后约 12 天退化，除非
人绒毛膜促性腺激素（hCG）加以挽救；若无它，孕酮下降，子宫
内膜脱落，胚胎随月经一并丢失。
\textbf{11.} $\log_2 10^{5} = 16.6$ 次倍增；到第 10 周时，
\omterm{def:b2:mammal-reproduction:placenta}{胎盘}已能自行分泌孕酮，\omterm{def:b2:mammal-reproduction:ovary}{黄体}
不再需要。
\textbf{12.} 卵子在前期 I 停滞数十年，连接二价体的
黏连蛋白降解，因此\omterm{def:b2:meiosis-heredity:meiosis}{减数分裂} I 的不分离随母亲年龄
急剧上升；精子则在两个月内重新生成。
\textbf{13.} \omterm{def:b2:mammal-reproduction:implantation}{滋养层}形成之后（第 5 天），羊膜形成之前（第
8 天）：一个\omterm{def:b2:mammal-reproduction:placenta}{胎盘}，两个囊。
\textbf{14.} $3\times 10^{-8}\times 1.2\times 10^{5}\times 20/3.5\times
10^{-4} = \qty{206}{mL/min}$.
\textbf{15.} $7\times 3.5 = \qty{24.5}{mL/min}$：余量为 8 倍。
\textbf{16.} $500\times 0.16 = \qty{80}{mL/min}$；提取率
$24.5/80 = \qty{31}{\%}$.
\textbf{17.} $5\times 3.5\times 1440 = \qty{25}{g/d}$；在 \qty{28}{L}
母体血液中；每日血流量 \qty{720}{L}：胎儿约摄取
流经的葡萄糖的 \qty{4}{\%}。
\textbf{18.} 胎儿血红蛋白在
\qty{30}{mmHg} 时饱和度为 \qty{77}{\%}；胎儿红细胞多，心输出量高，
其组织适于在低压下工作——它实际上
生活在珠穆朗玛峰顶。
\textbf{19.} 肺尚未工作，气体交换、营养与
排泄无需胎儿的肠或肾，维持妊娠的激素
亦由\omterm{def:b2:mammal-reproduction:placenta}{胎盘}分泌；但它无法阻挡所有毒物——酒精、
尼古丁及某些病毒仍能通过。
\textbf{20.} 宫颈扩张 $\to$ 下丘脑 $\to$ \omterm{prop:b2:mammal-reproduction:birth}{催产素} $\to$
收缩 $\to$ 更多扩张，直至分娩解除刺激。
足月前，孕酮占主导的子宫中\omterm{prop:b2:mammal-reproduction:birth}{催产素}受体很少，
且无缝隙连接，收缩不能扩散，使
正反馈回路无法闭合。
\textbf{21.} $750\times 2.9 = \qty{2.2}{MJ}$；成本 $2.2/0.8 =
\qty{2.7}{MJ}$.
\textbf{22.} 生长 $30\times 6 = \qty{0.18}{MJ}$，维持
$300\times 4 = \qty{1.2}{MJ}$：共 \qty{1.4}{MJ}，约占
乳汁能量的三分之二；其余用于活动、保温及损失。
\textbf{23.} $270\times 2.7 = \qty{730}{MJ}$，约为
妊娠期的两倍半。
\textbf{24.} 哺乳会升高\omterm{prop:b2:mammal-reproduction:lactation}{催乳素}，并抑制促使\omterm{def:b2:mammal-reproduction:ovary}{排卵}的激素的
脉冲式释放，因此\omterm{def:b2:mammal-reproduction:ovary}{排卵}要数月后才恢复；
若不避孕，分娩间隔因此可达
两至三年。
\textbf{25.} $P = 0.95$；受精后第 17 天检测呈阳性
（周期第 31 天）；氧气余量为 8 倍；一日的乳汁
使母体付出 \qty{2.7}{MJ}。
\end{solution}
